\documentclass[unicode,10pt,a4paper,oneside,numbers=endperiod,openany]{scrartcl}
\usepackage[backend=biber,style=numeric]{biblatex}
\addbibresource{references.bib}
\usepackage{tikz}
\usepackage{pgfplots}
\usepgfplotslibrary{colorbrewer} % better colors
\pgfplotsset{compat=1.18}
\usepackage{amsmath}
\usepackage{url}
\usepackage{xcolor}
\usepackage{subcaption}
\usepackage{listings}
\usepackage{multirow}
\usepackage{graphicx}
\lstdefinestyle{grayverbatim}{
  basicstyle=\ttfamily,%
  backgroundcolor=\color{gray!10},%
  frame=single,%
  framesep=3pt,%
  breaklines=true,%
  columns=fullflexible%
  numberstyle = \tiny,%
  stepnumber = 1,%
  numbersep = 5pt,%
  numbers = left%
}
\definecolor{lightbluehighlight}{RGB}{204,232,255}
\input{assignment.sty}

\begin{document}


\setassignment
\setduedate{\textbf{Friday} 29 May 2026, 18:00~(Europe/Zurich)}

\serieheader{High-Performance Computing Lab for CSE}{2026}
            {Student: Riccardo Polo}
            {Discussed with: FULL NAME}
            {Solution for Project 7}{}
\newline

\assignmentpolicy

\section{GPU architecture and performance characteristics [15 Points]}

\subsection{Peak performance [5 Points]}
\begin{lstlisting}[style=grayverbatim,caption={Output of the \texttt{deviceQuery} CUDA sample.},language=bash,label={lst:deviceQuery}]
ripolo@eu-login-23:~/cuda-samples/Samples/1_Utilities/deviceQuery/build$ srun --gpus=rtx_2080:1 ./deviceQuery
srun: job 66431453 queued and waiting for resources
srun: job 66431453 has been allocated resources
/cluster/home/ripolo/cuda-samples/Samples/1_Utilities/deviceQuery/build/./deviceQuery Starting...

 CUDA Device Query (Runtime API) version (CUDART static linking)

Detected 1 CUDA Capable device(s)

Device 0: "NVIDIA GeForce RTX 2080 Ti"
  CUDA Driver Version / Runtime Version          13.0 / 13.0
  CUDA Capability Major/Minor version number:    7.5
  Total amount of global memory:                 10821 MBytes (11346575360 bytes)
  (068) Multiprocessors, (064) CUDA Cores/MP:    4352 CUDA Cores
  GPU Max Clock rate:                            1545 MHz (1.54 GHz)
  Memory Clock rate:                             7000 Mhz
  Memory Bus Width:                              352-bit
  L2 Cache Size:                                 5767168 bytes
  Maximum Texture Dimension Size (x,y,z)         1D=(131072), 2D=(131072, 65536), 3D=(16384, 16384, 16384)
  Maximum Layered 1D Texture Size, (num) layers  1D=(32768), 2048 layers
  Maximum Layered 2D Texture Size, (num) layers  2D=(32768, 32768), 2048 layers
  Total amount of constant memory:               65536 bytes
  Total amount of shared memory per block:       49152 bytes
  Total shared memory per multiprocessor:        65536 bytes
  Total number of registers available per block: 65536
  Warp size:                                     32
  Maximum number of threads per multiprocessor:  1024
  Maximum number of threads per block:           1024
  Max dimension size of a thread block (x,y,z): (1024, 1024, 64)
  Max dimension size of a grid size    (x,y,z): (2147483647, 65535, 65535)
  Maximum memory pitch:                          2147483647 bytes
  Texture alignment:                             512 bytes
  Concurrent copy and kernel execution:          Yes with 3 copy engine(s)
  Run time limit on kernels:                     No
  Integrated GPU sharing Host Memory:            No
  Support host page-locked memory mapping:       Yes
  Alignment requirement for Surfaces:            Yes
  Device has ECC support:                        Disabled
  Device supports Unified Addressing (UVA):      Yes
  Device supports Managed Memory:                Yes
  Device supports Compute Preemption:            Yes
  Supports Cooperative Kernel Launch:            Yes
  Device PCI Domain ID / Bus ID / location ID:   0 / 34 / 0
  Compute Mode:
     < Default (multiple host threads can use ::cudaSetDevice() with device simultaneously) >

deviceQuery, CUDA Driver = CUDART, CUDA Driver Version = 13.0, CUDA Runtime Version = 13.0, NumDevs = 1
Result = PASS 
\end{lstlisting}

From listing \ref{lst:deviceQuery} we can derive that:
\begin{itemize}
  \item The number of Streaming Multiprocessors (SMs) is $n_{\text{SMs}} = 68$
  \item The number of CUDA cores per SM is $n_{\text{cores per SM}} = 64$ 
  \item The GPU maximumclock frequency is $f = 1.545$ GHz
\end{itemize}

From listing \ref{lst:deviceQuery} we can also see that the CUDA Compute Capability is $7.5$ and from reffff we get that
\begin{itemize}
  \item For single precision (FP32) FMAs $TP_{\text{FP32}}=\frac{64}{64}=1$
  \item For double precision (FP64) FMAs $TP_{\text{FP64}}=\frac{64}{2}=32$
\end{itemize}
Thus we can caompute the theoretical peak performance of an NVIDIA GeForce RTX 2080 Ti for single and double precision as:
\begin{align*}
  P_{\text{FP32}} &= n_{\text{SMs}} \cdot \frac{n_{\text{cores per SM}}}{TP_{\text{FP32}}} \cdot n_{FMA} \cdot f\\ &= 68 \cdot \frac{64}{1} \cdot 2 \text{FP32} \cdot 1.545 \text{GHz} = 13450,56 \text{ GFLOPS} \\
  P_{\text{FP64}} &= n_{\text{SMs}} \cdot \frac{n_{\text{cores per SM}}}{TP_{\text{FP64}}} \cdot n_{FMA} \cdot f \\ &= 68 \cdot \frac{64}{32} \cdot 2 \text{FP64} \cdot 1.545 \text{GHz} = 420,33  \text{ GFLOPS}
\end{align*}

From Listing \ref{lst:deviceQuery} we can get the following informations
\begin{itemize}
  \item Global Memory : $10821$ MiB $ =10.57 $ GiB
  \item L2 Cache Size: $5.5$ MiB
  \item L1 and Shared Memory per SM 96 KiB 
  \item Read only memory 64 Kibytes
  \item Registers 64 Ki registers per block, each holding 4 bytes (total 256 KB).
\end{itemize}
We obtain the memory bus width $W_{\text{mem}} = 352$ bits and memory clock frequency as $f_{\text{mem}} = 7$ GHz. So we can compute the (theoretical) peak memory bandwidth of a GPU as:
\begin{align*}
  B_{\text{GPU}} &= \frac{W_{\text{mem}}}{8} \cdot f_{\text{mem}} \cdot 2\\
  &= \frac{352}{8} \cdot 7 \text{GHz} \cdot 2 = 616 \text{ GB/s}
\end{align*}

\subsection{Performance model: A simple roofline model [5 Points]}
Compute ridge points
$$ I_{\text{ridge}} = \frac{P_{\text{peak}}}{B_{\text{GPU}}}$$
For single precision:
$$ I_{\text{ridge, FP32}} = \frac{13450.56 \text{ GFLOPS}}{616 \text{ GB/s}} = 21.83 \text{ GLOP/byte}$$
For double precision:
$$ I_{\text{ridge, FP64}} = \frac{420.33 \text{ GFLOPS}}{616 \text{ GB/s}} = 0.68 \text{ GLOP/byte}$$

We can retrieve a roofline plot for the NVIDIA GeForce RTX 2080 Ti remembering that $P_{\text{attainable}} = \min{P_{\text{peak}}, B_{\text{GPU}} \cdot I_{\text{ridge}}}$

\section{Introduction to CUDA [50 Points]}

\subsection{Implement the simple kernels [30 Points]}
For the double format we have following output:
\begin{lstlisting}[style=grayverbatim,caption={Output of the simple kernels for the double format.},language=bash,label={lst:simple_kernels_double}]
ripolo@eu-login-11:~/hpc_projects/project07/02cuda/simple_kernels$ srun --gpus=rtx_2080:1 ./benchmark
srun: job 171788 queued and waiting for resources
srun: job 171788 has been allocated resources
[benchmark] axpy: warmup=5 runs=25 avg_time=0.363086 ms
Compute utilization = 46.2073 GFLOPS/s
Memory utilization = 554.487 GB/s
first = 4, last = 4, expected = 4
[benchmark] add scaled diff: warmup=5 runs=25 avg_time=0.477494 ms
Compute utilization = 52.704 GFLOPS/s
Memory utilization = 562.176 GB/s
first = 1.6, last = 1.6, expected = 1.6
[benchmark] scaled diff: warmup=5 runs=25 avg_time=0.363736 ms
Compute utilization = 46.1247 GFLOPS/s
Memory utilization = 553.496 GB/s
first = 0.6, last = 0.6, expected = 0.6
[benchmark] scale: warmup=5 runs=25 avg_time=0.247412 ms
Compute utilization = 33.9054 GFLOPS/s
Memory utilization = 542.486 GB/s
first = 0.1, last = 0.1, expected = 0.1
[benchmark] linear combination: warmup=5 runs=25 avg_time=0.363924 ms
Compute utilization = 69.1512 GFLOPS/s
Memory utilization = 553.21 GB/s
first = 1.5, last = 1.5, expected = 1.5
ripolo@eu-login-11:~/hpc_projects/project07/02cuda/simple_kernels$
\end{lstlisting}

For float format we have the following output:
\begin{lstlisting}[style=grayverbatim,caption={Output of the simple kernels for the float format.},language=bash,label={lst:simple_kernels_float}]
ripolo@eu-login-11:~/hpc_projects/project07/02cuda/simple_kernels$ srun --gpus=rtx_2080:1 ./benchmark
srun: job 172188 queued and waiting for resources
srun: job 172188 has been allocated resources
[benchmark] axpy: warmup=5 runs=25 avg_time=0.182811 ms
Compute utilization = 91.7736 GFLOPS/s
Memory utilization = 550.642 GB/s
first = 4, last = 4, expected = 4
[benchmark] add scaled diff: warmup=5 runs=25 avg_time=0.239703 ms
Compute utilization = 104.988 GFLOPS/s
Memory utilization = 559.933 GB/s
first = 1.6, last = 1.6, expected = 1.6
[benchmark] scaled diff: warmup=5 runs=25 avg_time=0.182762 ms
Compute utilization = 91.798 GFLOPS/s
Memory utilization = 550.788 GB/s
first = 0.6, last = 0.6, expected = 0.6
[benchmark] scale: warmup=5 runs=25 avg_time=0.125165 ms
Compute utilization = 67.0205 GFLOPS/s
Memory utilization = 536.164 GB/s
first = 0.1, last = 0.1, expected = 0.1
[benchmark] linear combination: warmup=5 runs=25 avg_time=0.182696 ms
Compute utilization = 137.747 GFLOPS/s
Memory utilization = 550.989 GB/s
first = 1.5, last = 1.5, expected = 1.5
\end{lstlisting}

\subsection{Implement the scalar product kernel [20 Points]}
\begin{lstlisting}[style=grayverbatim,caption={Output of the scalar product kernel for the double format.},language=bash,label={lst:scalar_product_double}]
ripolo@eu-login-20:~/hpc_projects/project07/02cuda/02_dot_product$ srun --gpus=rtx_2080:1 ./dot_prod
srun: job 191159 queued and waiting for resources
srun: job 191159 has been allocated resources
[benchmark] hpc_dot: warmup=5 runs=25 avg_time=27.3736 ms
Compute utilization = 1.22579 GFLOPS/s
Memory utilization = 9.80636 GB/s
Got: 16777216, expected: 16777216
\end{lstlisting}


\section{Non-linear PDE mini-app with CUDA: The Fisher equation [35 Points]}

output is 
\begin{lstlisting}[style=grayverbatim,caption={Output of the Fisher equation mini-app.},language=bash,label={lst:fisher}]
  ripolo@eu-login-20:~/hpc_projects/project07/03_fisher$ srun --gpus=rtx_2080:1 ./main 128 100 0.005
srun: job 192180 queued and waiting for resources
srun: job 192180 has been allocated resources
================================================================================
                      Welcome to mini-stencil!
version   :: CUDA
mesh      :: 128 * 128 dx = 0.00787402
time      :: 100 time steps from 0 .. 0.005
iteration :: CG 300, Newton 50, tolerance 1e-06
================================================================================
--------------------------------------------------------------------------------
simulation took 0.221755 seconds
1516 conjugate gradient iterations, at rate of 6836.37 iters/second
300 Newton iterations
--------------------------------------------------------------------------------
### 128, 100, 1516, 300, 0.221755 ###
Goodbye!
\end{lstlisting}


Using float the output cd
\end{document}
